\documentclass[10pt]{article}
\usepackage[T1]{fontenc}
\usepackage[utf8]{inputenc}
\usepackage[french]{babel}
\usepackage{amsmath,amssymb,mathrsfs,mathpazo}
\usepackage{geometry}
\geometry{a4paper,margin=2cm}
\usepackage{graphicx,array,tabularx,tabularray,colortbl,multirow,multicol,booktabs,enumitem,pifont,ulem,fancybox,diagbox,cancel}
\usepackage{pstricks,pstricks-add,pst-plot,pst-tree,pst-node,pst-text,pst-func,pst-math,pst-bspline,pst-3dplot,pst-eucl,pst-solides3d}
\let\pscalculate\undefined
\usepackage{pst-bezier}
\usepackage{tikz,pgfplots,tkz-tab}
\usetikzlibrary{arrows,arrows.meta,patterns,calc,positioning,trees,decorations.markings,shapes.geometric,shapes.misc,shapes.arrows,intersections,decorations.pathmorphing,decorations.pathreplacing,backgrounds,shadows,babel}
\pgfplotsset{compat=1.18}
\usepgfplotslibrary{fillbetween}
\usepackage[np]{numprint}
\usepackage{listings,makecell,pas-tableur,twoopt}
\let\barre\undefined
\usepackage[pstricks]{tablvar}
\definecolor{darkgreen}{RGB}{0,100,0}
\providecommand{\red}{\color{red}}
\providecommand{\blue}{\color{blue}}
\providecommand{\green}{\color{green}}
\providecommand{\black}{\color{black}}
\providecommand{\white}{\color{white}}
\providecommand{\gray}{\color{gray}}
\providecommand{\cyan}{\color{cyan}}
\providecommand{\magenta}{\color{magenta}}
\providecommand{\yellow}{\color{yellow}}
\providecommand{\violet}{\color{violet}}
\providecommand{\orange}{\color{orange}}
\providecommand{\euro}{\texteuro}
\providecommand{\EUR}{\texteuro}
\providecommand{\ieme}{\textsuperscript{e}}
\providecommand{\degres}{\ensuremath{^\circ}}
\providecommand{\degree}{\ensuremath{^\circ}}
\providecommand{\vv}[1]{\overrightarrow{#1}}
\providecommand{\Coord}[1]{#1}
\providecommand{\decosix}{}
\providecommand{\starredbullet}{\ensuremath{\bullet}}
\pagestyle{empty}
\setlength{\parindent}{0pt}
\usepackage{adjustbox}\DeclareUnicodeCharacter{2236}{:}\DeclareUnicodeCharacter{2190}{\ensuremath{\leftarrow}}
\providecommand{\R}{\mathbb{R}}
\providecommand{\N}{\mathbb{N}}
\providecommand{\Z}{\mathbb{Z}}
\providecommand{\Q}{\mathbb{Q}}
\providecommand{\D}{\mathbb{D}}
\providecommand{\C}{\mathbb{C}}
\providecommand{\e}{{\mathrm{e}}}
\providecommand{\vect}[1]{\overrightarrow{#1}}
\providecommand{\vectt}[1]{\overrightarrow{\mathrm{#1}}}
\providecommand{\vvt}[1]{\overrightarrow{\mathrm{#1}}}
\providecommand{\np}[2][]{{#2~#1}}
\providecommand{\vv}[1]{{\overrightarrow{#1}}}
\renewcommand{\d}{{\mathrm{d}}}
\providecommand{\degres}{{^\circ}}
\providecommand{\degre}{{^\circ}}
\providecommand{\blue}{\color{blue}}
\providecommand{\red}{\color{red}}
\providecommand{\enskip}{\quad}
\providecommand{\llbracket}{\lbrack\!\lbrack}
\providecommand{\rrbracket}{\rbrack\!\rbrack}
\providecommand{\euro}{{\text{€}}}
\providecommand{\no}{{\text{n°}}}
\providecommand{\interval}[2]{{[#1;#2]}}
\providecommand{\intervFO}[2]{{[#1;#2[}}
\providecommand{\dotfill}{\ldots}

\providecommand{\ds}{\displaystyle}
\providecommand{\textding}[1]{#1}
\providecommand{\Oij}{\(\mathrm{O};\vec{\imath},\vec{\jmath}\)}
\providecommand{\Oijk}{\(\mathrm{O};\vec{\imath},\vec{\jmath},\vec{k}\)}
\providecommand{\Ouv}{\(\mathrm{O};\vec{u},\vec{v}\)}
\providecommand{\intervFF}[2]{\ensuremath{[#1;#2]}}
\renewcommand{\euro}{\ifmmode\text{\texteuro}\else\texteuro\fi}
\DeclareUnicodeCharacter{00B0}{\ifmmode^{\circ}\else\textdegree\fi}
\DeclareUnicodeCharacter{20AC}{\ifmmode\text{\texteuro}\else\texteuro\fi}
% Énoncés : Éducation nationale. Transcription et corrigés : APMEP, Denis Vergès et contributeurs. Découpage et classement : CoopMaths, Sébastien Lozano, Éric Elter et contributeurs. Adaptation de la présentation : Maths Explorer.
% CC BY-SA : https://coopmaths.fr/www/about/
% https://coopmaths.fr/alea/static/bac/2021/tex/bac_2021_05_ameriquenord_5_cor.tex
\begin{document}
\noindent\textit{APMEP / CoopMaths. CC BY-SA. Présentation adaptée par Maths Explorer.}\par\medskip
\smallskip













\begin{center}
\psset{unit=1.25cm,arrowsize=2pt 3}
\begin{pspicture*}(-0.5,-0.8)(6.4,4.4)
\psgrid[gridlabels=0pt,subgriddiv=5,gridcolor=gray](0,-1)(9,5)
\psaxes[linewidth=1.25pt]{->}(0,0)(0,-0.8)(6.4,4.4)
\psplot[plotpoints=2000,linewidth=1.25pt,linecolor=red]{0.1}{6.4}{x ln 1 add 4 mul x div}
\psline[linewidth=1.25pt](0,4)(6.4,4)\uput[u](5.8,4){$T$}
\uput[u](5.8,1.9){\red $\mathcal{C}_f$}
\uput[ur](1,4){A}
\psplotTangent[arrows=<->]{1.649}{1}{x ln 1 add 4 mul x div}
\psline[ArrowInside=->]{->}(1.639,0)(1.639,3.639)(0,3.639)
\uput[ur](1.649,3.639){B}
\end{pspicture*}
\end{center}

\medskip

\begin{enumerate}
\item 
A$(1~;~4) \in \mathcal{C}_f$, donc  $f(1) = 4$ et la courbe $\mathcal{C}_f$ admet une tangente horizontale $T$ au point A(1~;~4); le coefficient directeur de cette tangente en ce point est nul ou encore le nombre dérivé est nul : $f'(1) = 0$.
\end{enumerate}



\[f(x) = \dfrac{a + b \ln x}{x} \,\, 
\text{où }\, a \text{ et}\, b \text{ sont deux nombres réels}.\]

\begin{enumerate}[resume]
\item 
$f$ est une fonction quotient de fonctions dérivables sur $]0~;~ +\infty[$, le dénominateur ne s'annulant pas et sur cet intervalle :

$f'(x) = \dfrac{\frac{b}{x} \times x - 1(a + b\ln x)}{x^2} = \dfrac{b - a - b\ln x}{x^2}$.

\item 
En utilisant les résultats du \textbf{1.} :

$f(1) = \dfrac{a + b \ln 1}{1} = 4 \iff a = 4$ ;

$f'(1) = \dfrac{b - 4 - b\ln 1}{1^2} = 0 \iff b - 4 = 0 \iff b = 4$.
\end{enumerate}

Dans la suite de l'exercice, on admet que la fonction $f$ est définie pour tout réel $x$ de l'intervalle $]0~;~ +\infty[$ par:

\[f(x) = \dfrac{4 + 4\ln x}{x}.\]

\begin{enumerate}[resume]
\item ~
$\bullet~~$On sait que $\displaystyle\lim_{x \to 0} \dfrac{\ln x}{x} = - \infty$, donc 
$\displaystyle\lim_{x \to 0} f(x) = - \infty$ ;

$\bullet~~$On a $f(x) = \dfrac{4}{x} + \dfrac{4\ln x}{x}$.

On a $\displaystyle\lim_{x \to + \infty}\dfrac{4}{x} = 0$ et on sait que $\displaystyle\lim_{x \to + \infty}\dfrac{4\ln x}{x} = 0$, donc par somme de limites : 

$\displaystyle\lim_{x \to + \infty}f(x) = 0$.
\item 
On a donc sur $]0~;~+ \infty[$, \, $f'(x) = \dfrac{-4\ln x}{x^2}$ qui a pour signe celui de $-4\ln x$.

On sait que sur $]0~;~1[$, \, $\ln x < 0$, donc $f'(x) > 0$ sur ]0~;~1[ ;

Par contre sur $]1~;~+ \infty[$, \, $\ln x > 0$, donc $f'(x) < 0$ sur $]1~;~+ \infty[$ ;

$f'(1) = 0$, donc le point de coordonnées (1~;~4) est le maximum de la fonction sur $]0~;~+ \infty[$.

La fonction $f$ est donc croissante sur ]0~;~1[ de $- \infty$ à 4, puis décroissante sur $[1~;~+ \infty[$ de 4 à 0 avec un maximum 4 pour $x = 1$.
\item 
$f'$étant une fonction quotient de fonctions dérivables sur $]0~;~+ \infty[$, le dénominateur ne s'annulant pas est dérivable et sur cet intervalle :

$f''(x) = \dfrac{-4 \times \frac{1}{x} \times x^2 - 2x \times (- 4\ln x)}{x^4} = \dfrac{- 4x + 8x\ln x}{x^4} = \dfrac{- 4 + 8\ln x}{x^3}.$

\item 
La courbe présente un point d'inflexion lorsque la dérivée seconde s'annule. Or :

$f''(x) = 0 \iff \dfrac{- 4 + 8\ln x}{x^3} = 0 \iff - 4 + 8\ln x = 0 \iff -1 + 2\ln x= 0 \iff \ln x= \dfrac{1}{2} \iff x = \text{e}^{\frac{1}{2}} \approx 1,649$ $\left(\text{ou}\, \sqrt{\text{e}} \right)$.

L'ordonnée de ce point unique d'inflexion est $f\left(\text{e}^{\frac{1}{2}}\right) = \dfrac{4 + 4\times \frac{1}{2}}{\text{e}^{\frac{1}{2}}} = \dfrac{6}{\text{e}^{\frac{1}{2}}} = \dfrac{6} {\sqrt{\text{e}}} \approx 3,639$.

Ce point d'inflexion et la tangente en ce point sont indiqués sur la figure ci-dessus.
\end{enumerate}

\begin{center}{\fontencoding{U}\fontfamily{futs}\fontsize{18}{18}\selectfont m} 


\end{center}

\end{document}
